(1) 二维晶格模型的连续低能正则化为
\[
H(\boldsymbol k)=A k_x\sigma_x+A k_y\sigma_y+(M-B|\boldsymbol k|^2)\sigma_z,
\]
\(A>0\)、实数 \(M,B\ne0\)，并规定实物理 \(k\) 空间在无穷远可紧化。求两条能带及可能的闭隙条件，说明二次质量正则项 \(B k^2\) 对整数拓扑不可缺少的原因。
(2) 占据负能量能带，采用Berry曲率约定 \(\Omega_{xy}=i(\langle\partial_{k_x}u_-|\partial_{k_y}u_-\rangle-\langle\partial_{k_y}u_-|\partial_{k_x}u_-\rangle)\)。从归一化Bloch矢量 \(\hat{\boldsymbol d}\) 推导 \(\Omega_{xy}(\boldsymbol k)\)，精确给出包含质量与正则项的分子。
(3) 将动量平面紧化成球面，通过轴对称Bloch矢量角度或直接积分计算下能带Chern数 \(C=(2\pi)^{-1}\int_{\mathbb R^2}\Omega_{xy}\,d^2k\)。分别讨论 \(M/B>0\) 和 \(M/B<0\)，解释何时能量闭隙允许拓扑改变，并检验改变 \(A\) 幅值不改变整数。
(4) 在缓变质量界面 \(M(y)\) 从 \(y\to-\infty\) 的负值变为 \(y\to+\infty\) 的正值，先取局部 \(B=0\) 的一阶Dirac边缘近似。将 \(k_y\to-i\partial_y\)，解出随 \(y\) 正规化的界面态、其 \(\sigma_x\) 极化和 \(E(k_x)\)，核对其手征方向与体拓扑两侧差异。
